\IfFileExists{stacks-project.cls}{%
\documentclass{stacks-project}
}{%
\documentclass{amsart}
}

% For dealing with references we use the comment environment
\usepackage{verbatim}
\newenvironment{reference}{\comment}{\endcomment}
%\newenvironment{reference}{}{}
\newenvironment{slogan}{\comment}{\endcomment}
\newenvironment{history}{\comment}{\endcomment}

% For commutative diagrams we use Xy-pic
\usepackage[all]{xy}

% We use 2cell for 2-commutative diagrams.
\xyoption{2cell}
\UseAllTwocells

% We use multicol for the list of chapters between chapters
\usepackage{multicol}

% This is generally recommended for better output
\usepackage{lmodern}
\usepackage[T1]{fontenc}

% For cross-file-references
\usepackage{xr-hyper}

% Package for hypertext links:
\usepackage{hyperref}

% For any local file, say "hello.tex" you want to link to please
% use \externaldocument[hello-]{hello}
\externaldocument[introduction-]{introduction}
\externaldocument[conventions-]{conventions}
\externaldocument[sets-]{sets}
\externaldocument[categories-]{categories}
\externaldocument[topology-]{topology}
\externaldocument[sheaves-]{sheaves}
\externaldocument[sites-]{sites}
\externaldocument[stacks-]{stacks}
\externaldocument[fields-]{fields}
\externaldocument[algebra-]{algebra}
\externaldocument[brauer-]{brauer}
\externaldocument[homology-]{homology}
\externaldocument[derived-]{derived}
\externaldocument[simplicial-]{simplicial}
\externaldocument[more-algebra-]{more-algebra}
\externaldocument[smoothing-]{smoothing}
\externaldocument[modules-]{modules}
\externaldocument[sites-modules-]{sites-modules}
\externaldocument[injectives-]{injectives}
\externaldocument[cohomology-]{cohomology}
\externaldocument[sites-cohomology-]{sites-cohomology}
\externaldocument[dga-]{dga}
\externaldocument[dpa-]{dpa}
\externaldocument[sdga-]{sdga}
\externaldocument[hypercovering-]{hypercovering}
\externaldocument[schemes-]{schemes}
\externaldocument[constructions-]{constructions}
\externaldocument[properties-]{properties}
\externaldocument[morphisms-]{morphisms}
\externaldocument[coherent-]{coherent}
\externaldocument[divisors-]{divisors}
\externaldocument[limits-]{limits}
\externaldocument[varieties-]{varieties}
\externaldocument[topologies-]{topologies}
\externaldocument[descent-]{descent}
\externaldocument[perfect-]{perfect}
\externaldocument[more-morphisms-]{more-morphisms}
\externaldocument[flat-]{flat}
\externaldocument[groupoids-]{groupoids}
\externaldocument[more-groupoids-]{more-groupoids}
\externaldocument[etale-]{etale}
\externaldocument[chow-]{chow}
\externaldocument[intersection-]{intersection}
\externaldocument[pic-]{pic}
\externaldocument[weil-]{weil}
\externaldocument[adequate-]{adequate}
\externaldocument[dualizing-]{dualizing}
\externaldocument[duality-]{duality}
\externaldocument[discriminant-]{discriminant}
\externaldocument[derham-]{derham}
\externaldocument[local-cohomology-]{local-cohomology}
\externaldocument[algebraization-]{algebraization}
\externaldocument[curves-]{curves}
\externaldocument[resolve-]{resolve}
\externaldocument[models-]{models}
\externaldocument[functors-]{functors}
\externaldocument[equiv-]{equiv}
\externaldocument[pione-]{pione}
\externaldocument[etale-cohomology-]{etale-cohomology}
\externaldocument[proetale-]{proetale}
\externaldocument[relative-cycles-]{relative-cycles}
\externaldocument[more-etale-]{more-etale}
\externaldocument[trace-]{trace}
\externaldocument[crystalline-]{crystalline}
\externaldocument[spaces-]{spaces}
\externaldocument[spaces-properties-]{spaces-properties}
\externaldocument[spaces-morphisms-]{spaces-morphisms}
\externaldocument[decent-spaces-]{decent-spaces}
\externaldocument[spaces-cohomology-]{spaces-cohomology}
\externaldocument[spaces-limits-]{spaces-limits}
\externaldocument[spaces-divisors-]{spaces-divisors}
\externaldocument[spaces-over-fields-]{spaces-over-fields}
\externaldocument[spaces-topologies-]{spaces-topologies}
\externaldocument[spaces-descent-]{spaces-descent}
\externaldocument[spaces-perfect-]{spaces-perfect}
\externaldocument[spaces-more-morphisms-]{spaces-more-morphisms}
\externaldocument[spaces-flat-]{spaces-flat}
\externaldocument[spaces-groupoids-]{spaces-groupoids}
\externaldocument[spaces-more-groupoids-]{spaces-more-groupoids}
\externaldocument[bootstrap-]{bootstrap}
\externaldocument[spaces-pushouts-]{spaces-pushouts}
\externaldocument[spaces-chow-]{spaces-chow}
\externaldocument[groupoids-quotients-]{groupoids-quotients}
\externaldocument[spaces-more-cohomology-]{spaces-more-cohomology}
\externaldocument[spaces-simplicial-]{spaces-simplicial}
\externaldocument[spaces-duality-]{spaces-duality}
\externaldocument[formal-spaces-]{formal-spaces}
\externaldocument[restricted-]{restricted}
\externaldocument[spaces-resolve-]{spaces-resolve}
\externaldocument[formal-defos-]{formal-defos}
\externaldocument[defos-]{defos}
\externaldocument[cotangent-]{cotangent}
\externaldocument[examples-defos-]{examples-defos}
\externaldocument[algebraic-]{algebraic}
\externaldocument[examples-stacks-]{examples-stacks}
\externaldocument[stacks-sheaves-]{stacks-sheaves}
\externaldocument[criteria-]{criteria}
\externaldocument[artin-]{artin}
\externaldocument[quot-]{quot}
\externaldocument[stacks-properties-]{stacks-properties}
\externaldocument[stacks-morphisms-]{stacks-morphisms}
\externaldocument[stacks-limits-]{stacks-limits}
\externaldocument[stacks-cohomology-]{stacks-cohomology}
\externaldocument[stacks-perfect-]{stacks-perfect}
\externaldocument[stacks-introduction-]{stacks-introduction}
\externaldocument[stacks-more-morphisms-]{stacks-more-morphisms}
\externaldocument[stacks-geometry-]{stacks-geometry}
\externaldocument[moduli-]{moduli}
\externaldocument[moduli-curves-]{moduli-curves}
\externaldocument[examples-]{examples}
\externaldocument[exercises-]{exercises}
\externaldocument[guide-]{guide}
\externaldocument[desirables-]{desirables}
\externaldocument[coding-]{coding}
\externaldocument[obsolete-]{obsolete}
\externaldocument[fdl-]{fdl}
\externaldocument[index-]{index}

% Theorem environments.
%
\theoremstyle{plain}
\newtheorem{theorem}[subsection]{Theorem}
\newtheorem{proposition}[subsection]{Proposition}
\newtheorem{lemma}[subsection]{Lemma}

\theoremstyle{definition}
\newtheorem{definition}[subsection]{Definition}
\newtheorem{example}[subsection]{Example}
\newtheorem{exercise}[subsection]{Exercise}
\newtheorem{situation}[subsection]{Situation}

\theoremstyle{remark}
\newtheorem{remark}[subsection]{Remark}
\newtheorem{remarks}[subsection]{Remarks}

\numberwithin{equation}{subsection}

% Macros
%
\def\lim{\mathop{\mathrm{lim}}\nolimits}
\def\colim{\mathop{\mathrm{colim}}\nolimits}
\def\Spec{\mathop{\mathrm{Spec}}}
\def\Hom{\mathop{\mathrm{Hom}}\nolimits}
\def\Ext{\mathop{\mathrm{Ext}}\nolimits}
\def\SheafHom{\mathop{\mathcal{H}\!\mathit{om}}\nolimits}
\def\SheafExt{\mathop{\mathcal{E}\!\mathit{xt}}\nolimits}
\def\Sch{\mathit{Sch}}
\def\Mor{\mathop{\mathrm{Mor}}\nolimits}
\def\Ob{\mathop{\mathrm{Ob}}\nolimits}
\def\Sh{\mathop{\mathit{Sh}}\nolimits}
\def\NL{\mathop{N\!L}\nolimits}
\def\CH{\mathop{\mathrm{CH}}\nolimits}
\def\proetale{{pro\text{-}\acute{e}tale}}
\def\etale{{\acute{e}tale}}
\def\QCoh{\mathit{QCoh}}
\def\Ker{\mathop{\mathrm{Ker}}}
\def\Im{\mathop{\mathrm{Im}}}
\def\Coker{\mathop{\mathrm{Coker}}}
\def\Coim{\mathop{\mathrm{Coim}}}

% Boxtimes
%
\DeclareMathSymbol{\boxtimes}{\mathbin}{AMSa}{"02}

%
% Macros for moduli stacks/spaces
%
\def\QCohstack{\mathcal{QC}\!\mathit{oh}}
\def\Cohstack{\mathcal{C}\!\mathit{oh}}
\def\Spacesstack{\mathcal{S}\!\mathit{paces}}
\def\Quotfunctor{\mathrm{Quot}}
\def\Hilbfunctor{\mathrm{Hilb}}
\def\Curvesstack{\mathcal{C}\!\mathit{urves}}
\def\Polarizedstack{\mathcal{P}\!\mathit{olarized}}
\def\Complexesstack{\mathcal{C}\!\mathit{omplexes}}
% \Pic is the operator that assigns to X its picard group, usage \Pic(X)
% \Picardstack_{X/B} denotes the Picard stack of X over B
% \Picardfunctor_{X/B} denotes the Picard functor of X over B
\def\Pic{\mathop{\mathrm{Pic}}\nolimits}
\def\Picardstack{\mathcal{P}\!\mathit{ic}}
\def\Picardfunctor{\mathrm{Pic}}
\def\Deformationcategory{\mathcal{D}\!\mathit{ef}}
\usepackage{cleveref}
\crefname{originallemma}{original lemma}{original lemmas}
\Crefname{originallemma}{Original lemma}{Original lemmas}
\makeatletter
\newlabel{algebra-lemma-local-flat-ff}{{00HR}{source}{}{sourcepointer-1}{}}
\newlabel{algebra-lemma-local-flat-ff@cref}{{[originallemma][0][]00HR}{[1][0][]source}{}{}{}}
\makeatother
\crefname{originallemma}{original lemma}{original lemmas}
\Crefname{originallemma}{Original lemma}{Original lemmas}
\makeatletter
\newlabel{lemma-refine-descent-datum}{{023C}{source}{}{sourcepointer-2}{}}
\newlabel{lemma-refine-descent-datum@cref}{{[originallemma][0][]023C}{[1][0][]source}{}{}{}}
\makeatother
\crefname{originallemma}{original lemma}{original lemmas}
\Crefname{originallemma}{Original lemma}{Original lemmas}
\makeatletter
\newlabel{lemma-standard-fpqc-covering}{{023S}{source}{}{sourcepointer-3}{}}
\newlabel{lemma-standard-fpqc-covering@cref}{{[originallemma][0][]023S}{[1][0][]source}{}{}{}}
\makeatother
\crefname{originallemma}{original lemma}{original lemmas}
\Crefname{originallemma}{Original lemma}{Original lemmas}
\makeatletter
\newlabel{lemma-zariski-descent-effective}{{023E}{source}{}{sourcepointer-4}{}}
\newlabel{lemma-zariski-descent-effective@cref}{{[originallemma][0][]023E}{[1][0][]source}{}{}{}}
\makeatother
\crefname{originaldefinition}{original definition}{original definitions}
\Crefname{originaldefinition}{Original definition}{Original definitions}
\makeatletter
\newlabel{sites-definition-morphism-coverings}{{00VT}{source}{}{sourcepointer-5}{}}
\newlabel{sites-definition-morphism-coverings@cref}{{[originaldefinition][0][]00VT}{[1][0][]source}{}{}{}}
\makeatother
\crefname{originaldefinition}{original definition}{original definitions}
\Crefname{originaldefinition}{Original definition}{Original definitions}
\makeatletter
\newlabel{topologies-definition-fpqc-covering}{{022B}{source}{}{sourcepointer-6}{}}
\newlabel{topologies-definition-fpqc-covering@cref}{{[originaldefinition][0][]022B}{[1][0][]source}{}{}{}}
\makeatother
\crefname{originallemma}{original lemma}{original lemmas}
\Crefname{originallemma}{Original lemma}{Original lemmas}
\makeatletter
\newlabel{topologies-lemma-fpqc}{{022D}{source}{}{sourcepointer-7}{}}
\newlabel{topologies-lemma-fpqc@cref}{{[originallemma][0][]022D}{[1][0][]source}{}{}{}}
\makeatother
\title{Effectivity and full faithfulness of fpqc descent for quasi-coherent sheaves}
\author{The Stacks Project Authors}
\date{stacks-project:descent.tex:local-snapshot-20261012; GNU FDL 1.2 or later}
\begin{document}
\maketitle
\noindent\textit{Editorial note.} All delivered mathematical text is the author's own, copied from the pinned Stacks Project chapter file descent.tex; no statement or proof step was rewritten, reordered or removed.\par\medskip
\noindent\textit{Editorial note.} The proof ends with the author's own sentence that the final verification, namely that the canonical descent datum of the glued sheaf is isomorphic to the one started with, is omitted. That sentence is reproduced verbatim; the main line of the argument, faithfulness, fullness and essential surjectivity, is complete inside the counted span, so the omission is the closing routine verification rather than a missing block of the argument.\par\medskip
\noindent\textit{Editorial note.} The counted spans carry 7 navigation stubs; no bibliography citation occurs in the spans.\par\medskip
\noindent\textit{Editorial note.} The preamble span points at the chapter's own preamble.tex shipped with the frozen source, used verbatim as staged.\par\medskip
\noindent\textit{Allowed standard prerequisites (editorial).}\begin{itemize}
\item Standard affine fpqc coverings and effectivity of descent for them. Same chapter, Lemma lemma-standard-fpqc-covering, invoked at proof lines 1869-1871 and 1887-1888. Supplies effectivity and full faithfulness in the standard affine case.
\item Effectivity of Zariski descent for quasi-coherent sheaves. Same chapter, Lemma lemma-zariski-descent-effective, invoked at proof lines 1901-1903. Supplies the gluing of the local quasi-coherent modules over an affine open covering.
\item Pullback of descent data along a refinement. Same chapter, Lemma lemma-refine-descent-datum, invoked at proof lines 1866-1869. Supplies the comparison of descent data across refinements.
\item A local flat ring map with nilpotent-free residue extension is faithfully flat. Algebra, Lemma algebra-lemma-local-flat-ff, invoked at proof lines 1848-1851. Supplies the stalkwise conclusion for faithfulness.
\end{itemize}\medskip
\ifdefined\hypertarget\hypertarget{sourcepointer-1}{}\fi
\noindent\textit{Source reference.} Original lemma 00HR: sources/stacks-chow/tags.txt, source lines 653--653. Pointer outside the counted span, invoked once for context; the load-bearing argument of the counted statement is the author's proof reproduced in full in the proof span.\par\medskip
\ifdefined\hypertarget\hypertarget{sourcepointer-2}{}\fi
\noindent\textit{Source reference.} Original lemma 023C: sources/stacks-descent/pinned-original-source.tex, source lines 325--325. Pointer outside the counted span, invoked once for context; the load-bearing argument of the counted statement is the author's proof reproduced in full in the proof span.\par\medskip
\ifdefined\hypertarget\hypertarget{sourcepointer-3}{}\fi
\noindent\textit{Source reference.} Original lemma 023S: sources/stacks-descent/pinned-original-source.tex, source lines 1767--1767. Pointer outside the counted span, invoked once for context; the load-bearing argument of the counted statement is the author's proof reproduced in full in the proof span.\par\medskip
\ifdefined\hypertarget\hypertarget{sourcepointer-4}{}\fi
\noindent\textit{Source reference.} Original lemma 023E: sources/stacks-descent/pinned-original-source.tex, source lines 393--393. Pointer outside the counted span, invoked once for context; the load-bearing argument of the counted statement is the author's proof reproduced in full in the proof span.\par\medskip
\ifdefined\hypertarget\hypertarget{sourcepointer-5}{}\fi
\noindent\textit{Source reference.} Original definition 00VT: sources/stacks-chow/tags.txt, source lines 1159--1159. Pointer outside the counted span, invoked once for context; the load-bearing argument of the counted statement is the author's proof reproduced in full in the proof span.\par\medskip
\ifdefined\hypertarget\hypertarget{sourcepointer-6}{}\fi
\noindent\textit{Source reference.} Original definition 022B: sources/stacks-chow/tags.txt, source lines 2700--2700. Pointer outside the counted span, invoked once for context; the load-bearing argument of the counted statement is the author's proof reproduced in full in the proof span.\par\medskip
\ifdefined\hypertarget\hypertarget{sourcepointer-7}{}\fi
\noindent\textit{Source reference.} Original lemma 022D: sources/stacks-chow/tags.txt, source lines 2702--2702. Pointer outside the counted span, invoked once for context; the load-bearing argument of the counted statement is the author's proof reproduced in full in the proof span.\par\medskip
\begin{proposition}
\label{proposition-fpqc-descent-quasi-coherent}
Let $S$ be a scheme.
Let $\mathcal{U} = \{\varphi_i : U_i \to S\}$ be an fpqc covering, see
Topologies, Definition \ref{topologies-definition-fpqc-covering}.
Any descent datum on quasi-coherent sheaves
for $\mathcal{U} = \{U_i \to S\}$ is effective.
Moreover, the functor from the category of
quasi-coherent $\mathcal{O}_S$-modules to the category
of descent data with respect to $\mathcal{U}$ is fully faithful.
\end{proposition}

\begin{proof}
Let $S = \bigcup_{j \in J} V_j$ be an affine open covering.
For $j, j' \in J$ we denote $V_{jj'} = V_j \cap V_{j'}$ the intersection
(which need not be affine). For $V \subset S$ open we denote
$\mathcal{U}_V = \{V \times_S U_i \to V\}_{i \in I}$ which is a
fpqc-covering (Topologies, Lemma \ref{topologies-lemma-fpqc}).
By definition of an fpqc covering, we can find for each $j \in J$ a
finite set $K_j$, a map $\underline{i} : K_j \to I$,
affine opens $U_{\underline{i}(k), k} \subset U_{\underline{i}(k)}$,
$k \in K_j$ such that
$\mathcal{V}_j = \{U_{\underline{i}(k), k} \to V_j\}_{k \in K_j}$ is
a standard fpqc covering of $V_j$. And of course, $\mathcal{V}_j$
is a refinement of $\mathcal{U}_{V_j}$. Picture
$$
\xymatrix{
\mathcal{V}_j \ar[r] \ar@{~>}[d] &
\mathcal{U}_{V_j} \ar[r] \ar@{~>}[d] &
\mathcal{U} \ar@{~>}[d] \\
V_j \ar@{=}[r] & V_j \ar[r] & S
}
$$
where the top horizontal arrows are morphisms of families of
morphisms with fixed target (see
Sites, Definition \ref{sites-definition-morphism-coverings}).

\medskip\noindent
To prove the proposition you show successively the
faithfulness, fullness, and essential surjectivity of the
functor from quasi-coherent sheaves to descent data.

\medskip\noindent
Faithfulness. Let $\mathcal{F}$, $\mathcal{G}$ be quasi-coherent
sheaves on $S$ and let $a, b : \mathcal{F} \to \mathcal{G}$ be
homomorphisms of $\mathcal{O}_S$-modules.
Suppose $\varphi_i^*(a) = \varphi_i^*(b)$ for all $i$.
Pick $s \in S$. Then $s = \varphi_i(u)$ for some $i \in I$ and
$u \in U_i$. Since $\mathcal{O}_{S, s} \to \mathcal{O}_{U_i, u}$
is flat, hence faithfully flat
(Algebra, Lemma \ref{algebra-lemma-local-flat-ff}) we see
that $a_s = b_s : \mathcal{F}_s \to \mathcal{G}_s$. Hence $a = b$.

\medskip\noindent
Fully faithfulness. Let $\mathcal{F}$, $\mathcal{G}$ be quasi-coherent
sheaves on $S$ and let
$a_i : \varphi_i^*\mathcal{F} \to \varphi_i^*\mathcal{G}$ be
homomorphisms of $\mathcal{O}_{U_i}$-modules such that
$\text{pr}_0^*a_i = \text{pr}_1^*a_j$ on $U_i \times_U U_j$.
We can pull back these morphisms to get morphisms
$$
a_k :
\varphi_{i(k)}^*\mathcal{F}|_{U_{\underline{i}(k), k}}
\longrightarrow
\varphi_{i(k)}^*\mathcal{G}|_{U_{\underline{i}(k), k}}
$$
$k \in K_j$ with notation as above. Moreover,
Lemma \ref{lemma-refine-descent-datum} assures us
that these define a morphism between (canonical) descent data on
$\mathcal{V}_j$. Hence, by
Lemma \ref{lemma-standard-fpqc-covering}, we get correspondingly
unique morphisms $a_j : \mathcal{F}|_{V_j} \to \mathcal{G}|_{V_j}$.
To see that $a_j|_{V_{jj'}} = a_{j'}|_{V_{jj'}}$ we use that
both $a_j$ and $a_{j'}$ agree with the pullback of the morphism
$(a_i)_{i \in I}$ of (canonical) descent data to any covering
refining both $\mathcal{V}_{j, V_{jj'}}$ and
$\mathcal{V}_{j', V_{jj'}}$, and using the faithfulness already
shown. For example the covering
$\mathcal{V}_{jj'} =
\{V_k \times_S V_{k'} \to V_{jj'}\}_{k \in K_j, k' \in K_{j'}}$
will do.

\medskip\noindent
Essential surjectivity. Let $\xi = (\mathcal{F}_i, \varphi_{ii'})$
be a descent datum for quasi-coherent sheaves relative to the covering
$\mathcal{U}$. Pull back this descent datum to get descent data
$\xi_j$ for quasi-coherent sheaves relative to the coverings
$\mathcal{V}_j$ of $V_j$. By Lemma \ref{lemma-standard-fpqc-covering}
once again there exist
quasi-coherent sheaves $\mathcal{F}_j$ on $V_j$ whose associated
canonical descent datum is isomorphic to $\xi_j$. By fully faithfulness
(proved above) we see there are isomorphisms
$$
\phi_{jj'} :
\mathcal{F}_j|_{V_{jj'}}
\longrightarrow
\mathcal{F}_{j'}|_{V_{jj'}}
$$
corresponding to the isomorphism of descent data between the pullback
of $\xi_j$ and $\xi_{j'}$ to $\mathcal{V}_{jj'}$. To see that these
maps $\phi_{jj'}$ satisfy the cocycle condition we use faithfulness
(proved above) over the triple intersections $V_{jj'j''}$. Hence, by
Lemma \ref{lemma-zariski-descent-effective}
we see that the sheaves $\mathcal{F}_j$
glue to a quasi-coherent sheaf $\mathcal{F}$ as desired.
We still have to verify that the canonical descent datum relative to
$\mathcal{U}$ associated to $\mathcal{F}$ is isomorphic to the descent
datum we started out with. This verification is omitted.
\end{proof}

\end{document}
